\documentclass[10pt,a4paper]{amsart}
\usepackage[margin=19mm]{geometry}
\usepackage{amsmath,amssymb,mathtools}
\usepackage[scr=rsfs,cal=euler]{mathalfa}
\usepackage{xfrac}
\usepackage[unicode,hidelinks,bookmarks=false]{hyperref}
\newcommand{\ie}{i.e.\ }
\newcommand{\cf}{cf.\ }
\newcommand{\resp}{respectively\ }
\newcommand{\Subsection}{\subsection}
\providecommand{\nobreakdash}{\nobreak\hbox{-}}
\setlength{\emergencystretch}{2em}
\multlinegap0pt
\usepackage{bm}
\usepackage{bbm}
\usepackage{dsfont}
\usepackage{xspace}
\usepackage{cancel}
\usepackage{mathtools}
\usepackage{amsthm}
\makeatletter
\@ifundefined{lemma}{\newtheorem{lemma}{Lemma}}{}
\makeatother
\title[Decompositions of Tetrahedra into Height Functions]{Decompositions of Tetrahedra into Height Functions}
\author{Connor Castillo, Jackson Smith}
\date{Source: arXiv:2608.11217v1 (CC BY 4.0)}
\begin{document}
\maketitle

\noindent\emph{Source and licence.} Decompositions of Tetrahedra into Height Functions. Version arXiv:2608.11217v1 (CC BY 4.0), distributed under the Creative Commons Attribution 4.0 International licence, as recorded in the arXiv licence metadata for this exact version and shown on the abstract page https://arxiv.org/abs/2608.11217v1. Full rights notice: licenses/arxiv-2608.11217-CC-BY-4.0.txt.\par\smallskip

\noindent\textit{Editorial scope.} Counted unit: the Lemma whose source label is \texttt{L1} in the article (source0.tex lines 500--503, source label \texttt{L1}), delivered together with the article's own complete original proof of it (source0.tex lines 506--611, one \texttt{proof} environment of 106 lines). No mathematics is written, repaired or re-derived: statement and proof are the article's own text line for line. Required context retained uncounted: none. Every definition, display and label of those lines is retained unchanged. Omitted from this package and not redistributed: 1--499, 504--505, 612--1023. Nothing inside a retained excerpt is omitted or altered: every retained body is delivered exactly as the article's lines are written, so no omission receipt of any kind accompanies this package. Citations: the retained excerpts carry no citation command at all, so no bibliography entry of the article is transcribed and no reference is redistributed. Reference adaptation: none was needed and none was made; every reference inside the delivered unit resolves inside it, so no unresolved reference or citation is produced. Printed slips kept literal: no printed word, symbol, hypothesis or number of the counted statement or of its proof was altered. Layout and class: the article's own class is \texttt{mainresearch} with packages including ; the wrapper is the lane's pinned \texttt{amsart} editorial wrapper at 10pt on A4, which restates the theorem-like environments the retained text needs and adds the article's own macro definitions that the retained text needs (none), with extra wrapper packages (bm, bbm, dsfont, xspace, cancel, mathtools). The delivered numbering is the unit's own. Assets: no retained span contains a figure environment, a graphics-inclusion command, a PDF-page inclusion, a table or a TikZ picture; no image file is copied into this package, the package declares no asset dependency and no image file is redistributed with it. Licence: version arXiv:2608.11217v1 of this article is distributed under the Creative Commons Attribution 4.0 International licence, as recorded in the arXiv licence metadata for this exact version and shown on the abstract page https://arxiv.org/abs/2608.11217v1. A full rights notice is delivered at licenses/arxiv-2608.11217-CC-BY-4.0.txt. Characters: every retained excerpt is pure ASCII, so the delivered engine, XeLaTeX with its default Latin Modern text font, reports no missing character.
\begin{lemma}[Height-Dihedral Criterion]\label{L1}
	Let $T  \subset \mathbb{R}^3$ be a tetrahedron and $F$ be a face of $T$ with edges $e_1,e_2,e_3$.
	Then $T$ is a height function relative to the foundation $F$ if and only if the dihedral angles along the edges of $F$ are all not obtuse.
\end{lemma}

\begin{proof}
	$(\Rightarrow)$ Let $T$ be a height function relative to the foundation $F$. After possibly relabeling the vertices of $T$ we may assume $F = \operatorname{conv}\{v_0,v_1,v_2\}$. Notice that applying isometries and scalings to $\mathbb{R}^3$ does not alter the dihedral angles of $T$, as such, we will translate, rotate, reflect, and scale $T$ in order to produce a simpler environment for studying $\alpha_T([v_0,v_1]), \alpha_T([v_0,v_2])$ and $\alpha_T([v_1,v_2])$. After applying a translation, assume $v_0 = 0 \in \mathbb{R}^3$. Next, we apply a rigid rotation of $\mathbb{R}^3$ such that the affine hull of $F$ coincides with the $xy$-plane and the edge $[v_0,v_1]$ lies along the positive $x$-axis. In particular, this means that $v_1,v_2 \in \mathbb{R}^3$ have zero $z$ component. Hence, we may write
	\begin{equation*}
		v_1 = (x_1,0,0), \quad v_2 = (x_2,y_2,0)
	\end{equation*}
	with $x_1>0$. If necessary, reflect $T$ across the $xz$ plane to ensure $y_2>0$, and scale $T$ so that
	\begin{equation*}
		v_3 = (x_3,y_3,1)\quad
	\end{equation*}
	
	Considering the face $K = \operatorname{conv}\{v_0,v_1,v_3\}$ we move to demonstrate that $\theta_1$, the dihedral angle that occurs at the edge $[v_0,v_1]$ shared by $F$ and $K$ must be non-obtuse. Let $P_1 = \operatorname{aff}\{v_0,v_1,v_2\}$ be the plane containing $F$, and $P_2 = \operatorname{aff}\{v_0,v_1,v_3\}$ be the plane containing $K$. Since $v_0=0$ these planes simplify as,
	\[
	P_1= \operatorname{span} \{v_1,v_2\} \quad \text{ and } \quad P_2= \operatorname{span} \{v_1,v_3\}
	\]
	Because $v_1,v_2$ are linearly independent vectors in the $xy$ plane we conclude that $P_1$ is the $xy$ plane, as such, the vector $n=(0,0,1)$ is normal to $P_1$. Also, 
	\begin{equation*}
		m = v_1\times v_3 = (x_1,0,0) \times (x_3,y_3,1) = (0,-x_1,x_1y_3)
	\end{equation*}
	is normal to $P_2$. Then the dot product between both normal vectors is 
	\begin{align*}
		\langle m , n \rangle  &= (0,-x_1,x_1y_3) \cdot (0,0,1) \\
		&= x_1y_3
	\end{align*}
	Because $x_1>0$, it remains to show that $y_3 \ge 0$. Since $F$ is a foundation of $T$,
	\begin{equation*}
		\pi_F(v_3)= (x_3,y_3,0) \in F
	\end{equation*}
	Then by the convexity of $F$, there exists $r,s,t \ge 0$ with $r+s+t = 1$ such that
	\begin{equation*}
		(x_3,y_3,0) = rv_0+sv_1+tv_2 = sv_1+tv_2
	\end{equation*}
	Hence 
	\begin{equation*}
		(x_3,y_3,0) = (sx_1+tx_2,ty_2,0) \implies y_3 =ty_2
	\end{equation*}
	Since $y_2 >0$ and $t\ge 0$, we get that $y_3 \ge 0$, and therefore
	\begin{equation*}
		\langle n , m \rangle = x_1y_3 \ge 0 
	\end{equation*}
	Due to the dot product angle formula, we have that $\cos(\alpha_T(e_1)) \ge 0 $, and subsequently $\alpha_T([v_0,v_1]) \le \frac{\pi}{2}$. Because the initial normalization and coordinate transformations may be applied relative to any edge of $F$, the preceding argument applies identically to the remaining edges of $F$. Therefore, the remaining angles $\alpha_T([v_0,v_2])$ and $ \alpha_T([v_1,v_2])$ are non-obtuse.
	
	
	
	
	
	
	$(\Leftarrow)$ Suppose the dihedral angles of the edges of $F=\operatorname{conv}\{v_0,v_1,v_2\}$, given by, 
	\[
	\alpha_T([v_0,v_1]),\alpha_T([v_0,v_2]),\alpha_T([v_1,v_2])
	\]
	are all non-obtuse. To show $F$ is a foundation of $T$ we must show that $\pi_F(x) \in F$ for all $x\in T$. In fact, it is sufficient to show that $\pi_F(v_3)\in F$. Indeed, if $\pi_F(v_3) \in F = \operatorname{conv}\{v_0,v_1,v_2\}$ then, 
	\begin{equation*}
		\pi_F(v_3) = \sigma_0v_0+\sigma_1v_1+\sigma_2v_2 ,  \quad \text{ with } \sigma_i \ge 0 , \text{ and } \sum^2_{i=0} \sigma_i = 1
	\end{equation*}
	By the convexity of $T$, we may express any $x\in T$ as a convex combination:
	\begin{equation*}
		x = \lambda_0v_0+\lambda_1v_1+\lambda_2v_2+\lambda_3v_3 ,  \quad \text{ with } \lambda_i \ge 0 , \text{ and } \sum^3_{i=0} \lambda_i = 1
	\end{equation*}
	Then the orthogonal projection of $x$ onto the plane $\operatorname{aff}(F)$ is given by,
	\begin{align*}
		\pi_F(x) &= \lambda_0\pi_F(v_0)+\lambda_1\pi_F(v_1) +\lambda_2\pi_F(v_2)+\lambda_3 \pi_F(v_3) \\
		&= \lambda_0v_0+\lambda_1v_1+\lambda_2v_2 +\lambda_3 \pi_F(v_3)\\
		&= (\lambda_0 + \lambda_3\sigma_0 )v_0+(\lambda_1+\lambda_3\sigma_1)v_1+(\lambda_2 + \lambda_3\sigma_2) v_2 
	\end{align*}
	Observe that,
	\[
	\lambda_i+\lambda_3\sigma_i \ge 0 \quad \text{ and } \quad \sum_{i=0}^2(\lambda_i+\lambda_3\sigma_i) = \sum_{i=0}^3\lambda_i=1
	\]
	Therefore $\pi_F(x)\in \operatorname{conv}\{v_0,v_1,v_2\} = F$.
	
	With this in mind, we aim to show that $\pi_F(v_3) \in F$. Note that applying isometries and scaling $\mathbb{R}^3$ does not change whether or not $F\subset T$ is a foundation. As such, we will apply the exact same translation, rotations, reflections and scaling to $T$ as we did above.
	
	We start by extending the edge $[v_0,v_1]$ into a line $L=\mathrm{span}\{v_1\}$ (since $v_0 = 0$ and $v_1=(x_1,0,0)$, $L$ is simply the  $x$ axis). Observe that $L$ partitions the plane containing $F$ (the $xy$ plane) into two half spaces, namely
	\begin{equation*}
		H_1^+=\{(x,y):y\geq0\} \quad \text{and} \quad  H_1^-=\{(x,y):y<0\}
	\end{equation*}
	By our configuration of $T$, $v_2 = (x_2,y_2,0)$ with $y_2>0$. Hence, $v_2\in H_1^+$ and $F=\operatorname{conv}\{v_0,v_1,v_2\}\subset H_1^+$. 
	
	Note that the dihedral angle $\alpha_T([v_0,v_1])$ is given by the angle between the planes $\operatorname{aff}\{v_0,v_1,v_2\}$ and $\operatorname{aff}\{v_0,v_1,v_3\}$. Another way to express $\alpha_T([v_0,v_1])$ is as the angle between two vectors $n\in \operatorname{aff}\{v_0,v_1,v_2\}, m\in \operatorname{aff}\{v_0,v_1,v_3\}$ that are both orthogonal to the intersection line of the planes. It will soon be clear why we wish to express $\alpha_T([v_0,v_1])$ in this way, but first we aim to construct these vectors $n\in \operatorname{aff}\{v_0,v_1,v_2\}, m\in \operatorname{aff}\{v_0,v_1,v_3\}$ that are both orthogonal to $L=\operatorname{aff}\{v_0,v_1,v_2\} \cap \operatorname{aff}\{v_0,v_1,v_3\}$.
	
	Define $n$ as follows:
	\[
	n=(0,1,0)\in \operatorname{aff}\{v_0,v_1,v_2\}
	\]  
	Clearly, $n \perp L$ because $L$ is the $x$ axis.
	
	Next, let $\pi_L(v_3) = (x_3,0,0) \in L$, the orthogonal projection of $v_3$ onto the line $L$. By the properties of orthogonal projection,
	\begin{equation*}
		v_3-\pi_L(v_3) \perp L.
	\end{equation*}
	Moreover, $\pi_L(v_3)\in L$ means that $\pi_L(v_3) =tv_1$ for some $t\in \mathbb{R}$. Then, 
	\[
	v_3-\pi_L(v_3) = v_3-tv_1\in \operatorname{span}\{v_1,v_3\}=\operatorname{aff}\{v_0,v_1,v_3\}
	\]
	So, if we let $m= v_3-\pi_L(v_3)$ we have our desired vectors.
	Thus, $\theta_1$ is given by the angle between $n$ and $m$.
	Because $\theta_1$ is non-obtuse, it follows that,
	\begin{equation*}
		0\le n\cdot m = (0,1,0)\cdot(0,y_3,1) = y_3.
	\end{equation*}
	Therefore, $\pi_F(v_3) = (x_3,y_3,0) \in H_1^+$. Repeatedly applying this argument to the remaining     two edges of $F$ yields $\pi_F(v_3) \in H_1^+\cap H_2^+ \cap H_3^+ =F$. Therefore $F\subset T$ is a foundation and $T$ is a height function.
	
	
	
	
\end{proof}

\end{document}
